\chapter{Denoising prin prag}
\label{ch:denoising}

Zgomotul produce coeficienți wavelet mici. Semnalul util produce coeficienți
mari. Un prag care taie sau atenuează coeficienții mici, apoi reconstruiește,
elimină zgomotul fără să tocească muchiile, spre deosebire de un filtru
trece-jos clasic, care ar netezi și detaliile utile odată cu zgomotul.

Estimarea zgomotului folosește abaterea mediană absolută a subbenzii celei
mai fine de detaliu,
\[
\hat\sigma = \frac{\mathrm{median}(|d_{\text{fin}}|)}{0.6745},
\]
o alegere robustă: spre deosebire de o deviație standard calculată direct,
nu e trasă în sus de coeficienții mari care poartă semnal.
Pragul universal al lui Donoho și Johnstone,
$\lambda_{\text{universal}} = \hat\sigma\sqrt{2\ln N}$~\cite{donohojohnstone1994},
iar cele două reguli de trunchiere sunt
\[
\eta_{\text{soft}}(x,\lambda) = \mathrm{sign}(x)\max(|x|-\lambda,\,0), \qquad
\eta_{\text{hard}}(x,\lambda) = x\cdot \mathbb{1}(|x|>\lambda).
\]
Soft micșorează continuu coeficienții spre zero, fără discontinuități. Hard
elimină complet ce e sub prag și păstrează exact restul, cu riscul unor
mici artefacte la graniță (Figura~\ref{fig:denoise-theory}).

\begin{figure}[htbp]
\centering
\includegraphics[width=\textwidth]{41-denoise-theory.png}
\caption{Thresholding moale versus dur și distribuția tipică a coeficienților,
  cu pragul universal Donoho-Johnstone alături de pragul $2\hat\sigma$ folosit
  efectiv de demo.}
\label{fig:denoise-theory}
\end{figure}

Pragul folosit efectiv pe imagini e $2\hat\sigma$, mai mic decât pragul
universal din formulă. Abaterea e deliberată și documentată direct în cod. La
$\hat\sigma=15$ și $N=512^2$ (o imagine $512\times 512$), pragul universal
iese în jur de 68, atât de agresiv încât câștigul de raport semnal-zgomot
scade spre $0$~dB: teoria din spatele lui presupune un semnal neted plus
zgomot alb, o ipoteză prea optimistă pentru texturile unei fotografii.
Pragul $2\hat\sigma$ păstrează detaliul și obține în jur de $3$~dB câștig.
Demo-ul de ECG (Capitolul~\ref{ch:ecg}) folosește în schimb chiar pragul
universal, pe un semnal 1D sintetic, unde ipoteza semnal neted plus zgomot
alb se potrivește mult mai bine decât pe o imagine.

Pe o imagine reală (Figura~\ref{fig:denoise-demo}), același proces rulează
cu cifrele la vedere: dincolo de imaginea curățată apare și ce anume s-a
eliminat, amplificat de zece ori, zgomot granular pur, fără nicio structură
a imaginii originale, semn că a rămas doar semnalul.

\begin{figure}[htbp]
\centering
\includegraphics[width=\textwidth]{42-denoise-demo.png}
\caption{Denoising pe o imagine reală: reziduul amplificat $\times 10$ e
  zgomot granular pur, fără nicio structură a imaginii, semn că a rămas
  doar zgomotul.}
\label{fig:denoise-demo}
\end{figure}
